
%%% fhc mu-like

\begin{table}[htb]
\centering
\begin{tabular}{| c | c | c |}
\hline
  {\bf Source of uncertainty}    & \multicolumn{1}{c|}{{\bf ${\delta}N_{SK}$/$N_{SK}$}}\\
                                 & \multicolumn{1}{c|}{(Asimov A)}\\

\hline
\hline  SKDet+FSI+SI                   & 4.11\%                        \\
\hline
\hline
  Flux                           & 3.26\%                        \\
\hline
\hline
  2p-2h (corr)                     & 3.53\%                        \\
  2p-2h-bar (corr)                  & 0.20\%                        \\
  NC other (uncorr)              & 0.79\%                        \\
  NC 1gamma (uncorr)             & 0.00\%                        \\
  XSec nue/numu (uncorr)                 & 0.01\%                        \\
\hline
  XSec Tot (corr)                  & 3.96\%                        \\
  XSec Tot                       & 4.04\%                        \\
\hline
\hline
  Flux+XSec                      & 2.75\%                        \\
\hline
\hline
  Flux+XSec+SKDet+FSI+SI         & 5.03\%                        \\
\hline
\hline
  Oscillations                   & 0.03\%                        \\
\hline
\hline
  All                            & 5.00\%                        \\
\hline
\hline

\end{tabular}
\caption{
 Effect of 1$\sigma$ variations of the groups of systematics parameters on total number of
 1-ring muon-like events for neutrino mode obtained by performing 10k toys MC.
The systematic parameter values measured with the ND280 data 
shown in 
\cref{tab:banff_uncertainties_flux_fhc},
~\ref{tab:banff_uncertainties_flux_rhc}
~\ref{tab:banff_uncertainties_xsec}
and
~\ref{tab:banff_uncertainties_skdetfsi}
and the oscillation parameters of the Asimov data set A shown in \cref{tab:osc_params_joint_asimov}
are used.
Since the parameters $\theta_{23}$ and $\Delta m^2_{32}$ ($\Delta m^2_{13}$)
will be measured,
only the uncertainty on the parameters 
$\sin^2 \theta_{13}$,
$\sin^2 \theta_{12}$
and 
$\Delta m^2_{21}$ 
are taken into account.
The correction of the post-fit BANFF parameters as well as the additional covariance matrix 
used to take account for bias in the BANFF fit is not applied.
}                                                                                                                                                                                          
\label{tab:nsk_systematic_table_summary_joint_asimov_A_fhc_mulike_NoFluxBias}
\end{table}




%%%%% fhc e-like


\begin{table}[htb]
\centering
\begin{tabular}{| c | c | c |}
\hline
  {\bf Source of uncertainty}    & \multicolumn{1}{c|}{{\bf ${\delta}N_{SK}$/$N_{SK}$}}\\
                                 & \multicolumn{1}{c|}{(Asimov A)}\\

\hline
\hline  SKDet+FSI+SI                   & 3.47\%                        \\
\hline
  Flux                           & 3.40\%                        \\
\hline
\hline
  2p-2h (corr)                     & 3.95\%                        \\
  2p-2h-bar (corr)                  & 0.05\%                        \\
  NC other (uncorr)              & 0.16\%                        \\
  NC 1gamma (uncorr)             & 1.49\%                        \\
  XSec nue/numu (uncorr)                 & 2.63\%                        \\
\hline
  XSec Tot (corr)                  & 4.17\%                        \\
  XSec Tot                       & 5.15\%                        \\
\hline
\hline
  Flux+XSec                      & 4.08\%                        \\
\hline
  Flux+XSec+SKDet+FSI+SI         & 5.51\%                        \\
\hline
\hline
  Oscillations                   & 4.19\%                        \\
\hline
\hline
  All                            & 6.93\%                        \\
\hline
\hline

\end{tabular}
\caption{
 Effect of 1$\sigma$ variations of the groups of systematics parameters on total number of
 1-ring electron-like events for neutrino mode obtained by performing 10k toys MC.
The systematic parameter values measured with the ND280 data 
shown in 
\cref{tab:banff_uncertainties_flux_fhc},
~\ref{tab:banff_uncertainties_flux_rhc}
~\ref{tab:banff_uncertainties_xsec}
and
~\ref{tab:banff_uncertainties_skdetfsi}
and the oscillation parameters of the Asimov data set A shown in \cref{tab:osc_params_joint_asimov}
are used.
Since the parameters $\theta_{23}$ and $\Delta m^2_{32}$ ($\Delta m^2_{13}$)
will be measured,
only the uncertainty on the parameters 
$\sin^2 \theta_{13}$,
$\sin^2 \theta_{12}$
and 
$\Delta m^2_{21}$ 
are taken into account.
The correction of the post-fit BANFF parameters as well as the additional covariance matrix 
used to take account for bias in the BANFF fit is not applied.
}                                                                                                                                                                                          
\label{tab:nsk_systematic_table_summary_joint_asimov_A_fhc_elike_NoFluxBias}
\end{table}



%%%%%% rhc mu-like


\begin{table}[htb]
\centering
\begin{tabular}{| c | c | c |}
\hline
  {\bf Source of uncertainty}    & \multicolumn{1}{c|}{{\bf ${\delta}N_{SK}$/$N_{SK}$}}\\
                                 & \multicolumn{1}{c|}{(Asimov A)}\\

\hline
\hline  SKDet+FSI+SI                   & 3.94\%                        \\
\hline
\hline
  Flux                           & 3.40\%                        \\
\hline
\hline
  2p-2h (corr)                     & 3.00\%                        \\
  2p-2h-bar (corr)                  & 1.76\%                        \\
  NC other (uncorr)              & 0.76\%                        \\
  NC 1gamma (uncorr)             & 0.00\%                        \\
  XSec nue/numu (uncorr)                 & 0.00\%                        \\
\hline  
  XSec Tot (corr)                  & 4.20\%                        \\
  XSec Tot                       & 4.27\%                        \\
\hline
\hline
  Flux+XSec                      & 2.90\%                        \\
\hline
\hline
  Flux+XSec+SKDet+FSI+SI         & 4.90\%                        \\
\hline
\hline
  Oscillations                   & 0.03\%                        \\
\hline
\hline
  All                            & 4.91\%                        \\
\hline
\hline

\end{tabular}

\caption{
 Effect of 1$\sigma$ variations of the groups of systematics parameters on total number of
 1-ring muon-like events for antineutrino mode obtained by performing 10k toys MC.
The systematic parameter values measured with the ND280 data 
shown in 
\cref{tab:banff_uncertainties_flux_fhc},
~\ref{tab:banff_uncertainties_flux_rhc}
~\ref{tab:banff_uncertainties_xsec}
and
~\ref{tab:banff_uncertainties_skdetfsi}
and the oscillation parameters of the Asimov data set A shown in \cref{tab:osc_params_joint_asimov}
are used.
Since the parameters $\theta_{23}$ and $\Delta m^2_{32}$ ($\Delta m^2_{13}$)
will be measured,
only the uncertainty on the parameters 
$\sin^2 \theta_{13}$,
$\sin^2 \theta_{12}$
and 
$\Delta m^2_{21}$ 
are taken into account.
The correction of the post-fit BANFF parameters as well as the additional covariance matrix 
used to take account for bias in the BANFF fit is not applied.
}                                                                                                                                                                                          
\label{tab:nsk_systematic_table_summary_joint_asimov_A_rhc_mulike_NoFluxBias}
\end{table}



%%%%%% rhc e-like

\begin{table}[htb]
\centering
\begin{tabular}{| c | c | c |}
\hline
  {\bf Source of uncertainty}    & \multicolumn{1}{c|}{{\bf ${\delta}N_{SK}$/$N_{SK}$}}\\
                                 & \multicolumn{1}{c|}{(Asimov A)}\\

\hline
\hline  SKDet+FSI+SI                   & 3.94\%                        \\
\hline
\hline
  Flux                           & 3.53\%                        \\
\hline
\hline
  2p-2h (corr)                     & 3.11\%                        \\
  2p-2h-bar (corr)                  & 2.24\%                        \\
  NC other (uncorr)              & 0.33\%                        \\
  NC 1gamma (uncorr)             & 2.98\%                        \\
  XSec nue/numu (uncorr)                 & 1.52\%                        \\
\hline  
  XSec Tot (corr)                  & 4.44\%                        \\
  XSec Tot                       & 5.57\%                        \\
\hline
\hline
  Flux+XSec                      & 4.39\%                        \\
\hline
\hline
  Flux+XSec+SKDet+FSI+SI         & 6.11\%                        \\
\hline
\hline
  Oscillations                   & 3.98\%                        \\
\hline
\hline
  All                            & 7.24\%                        \\
\hline
\hline

\end{tabular}
\caption{
 Effect of 1$\sigma$ variations of the groups of systematics parameters on total number of
 1-ring electron-like events for antineutrino mode obtained by performing 10k toys MC.
The systematic parameter values measured with the ND280 data 
shown in 
\cref{tab:banff_uncertainties_flux_fhc},
~\ref{tab:banff_uncertainties_flux_rhc}
~\ref{tab:banff_uncertainties_xsec}
and
~\ref{tab:banff_uncertainties_skdetfsi}
and the oscillation parameters of the Asimov data set A shown in \cref{tab:osc_params_joint_asimov}
are used.
Since the parameters $\theta_{23}$ and $\Delta m^2_{32}$ ($\Delta m^2_{13}$)
will be measured,
only the uncertainty on the parameters 
$\sin^2 \theta_{13}$,
$\sin^2 \theta_{12}$
and 
$\Delta m^2_{21}$ 
are taken into account.
The correction of the post-fit BANFF parameters as well as the additional covariance matrix 
used to take account for bias in the BANFF fit is not applied.
}                                                                                                                                                                                          
\label{tab:nsk_systematic_table_summary_joint_asimov_A_rhc_elike_NoFluxBias}
\end{table}


















%% \begin{table}
%% \center
%% \begin{tabular}{| c || c | c || c | c |}
%% \hline
%% Source of uncertainty & $\mu$-like RMS/Mean &e-like RMS/Mean \\
%% \hline
%% SKDet & 0.007 & 0.016 \\ \hline
%% FSI+SI & 0.026 & 0.035 \\ \hline
%% SKDet+FSI+SI & 0.026 & 0.038 \\ \hline
%% \hline
%% Flux & 0.017 & 0.017 \\ \hline
%% XSec NC other (uncorr) & 0.000 & 0.002 \\ \hline
%% XSec NC $1\gamma$ (uncorr)  & 0.000  & 0.015 \\ \hline
%% XSec $\nu_e$ / $\nu_{\mu}$ (uncorr)  & 0.000  & 0.031 \\ \hline
%% XSec Tot  & 0.017  & 0.042 \\ \hline
%%  \hline
%% Flux+XSec  & 0.018  & 0.041 \\ \hline
%%  \hline
%% Flux+XSec+SDet+FSI+SI  & 0.032  & 0.057 \\ \hline
%%  \hline
%% Oscillations  & 0.000  & 0.008 \\ \hline
%%  \hline
%% All  & 0.032  & 0.058 \\ \hline
%% \end{tabular}
%% \caption{
%%  Effect of 1$\sigma$ variations of the several source of systematic uncertainties 
%%  on the ratio between the
%%  total number of events in neutrino and antineutrino modes
%%  for 1-ring electron- and muon-like samples
%%  obtained by performing 10k toys MC.
%%  The systematic parameter values measured with the ND280 data 
%% shown in 
%% tab.~\ref{tab:banff_uncertainties_flux_fhc},
%% ~\ref{tab:banff_uncertainties_flux_rhc}
%% ~\ref{tab:banff_uncertainties_xsec}
%% and
%% ~\ref{tab:banff_uncertainties_skdetfsi}
%% and the oscillation parameters of the Asimov data set A shown in tab.~\ref{tab:osc_params_joint_asimov}
%% are used.
%% Since the parameters $\theta_{23}$ and $\Delta m^2_{32}$ ($\Delta m^2_{13}$)
%% will be measured,
%% only the uncertainty on the parameters 
%% $\sin^2 \theta_{13}$,
%% $\sin^2 \theta_{12}$
%% and 
%% $\Delta m^2_{21}$ 
%% are taken into account.
%% }             
%% \label{tab:ratio_systematic_table_summary_asimov_A}               
%% \end{table}
